\documentclass[10pt,a4paper]{amsart}
\usepackage[margin=19mm]{geometry}
\usepackage{amsmath,amssymb,mathtools}
\usepackage[scr=rsfs,cal=euler]{mathalfa}
\usepackage{xfrac}
\usepackage[unicode,hidelinks,bookmarks=false]{hyperref}
\newcommand{\ie}{i.e.\ }
\newcommand{\cf}{cf.\ }
\newcommand{\resp}{respectively\ }
\newcommand{\Subsection}{\subsection}
\providecommand{\nobreakdash}{\nobreak\hbox{-}}
\setlength{\emergencystretch}{2em}
\multlinegap0pt
\usepackage{amssymb}
\usepackage{xcolor}
\usepackage[shortlabels]{enumitem}
\usepackage{float}
\newtheorem{lemma}{Lemma}
\newtheorem{proposition}{Proposition}
\def\CC{\mathbb{C}}
\def\RR{\mathbb{R}}
\def\S{\mathcal{S}}
\def\Sym{\operatorname{Sym}}
\def\b{\beta}
\def\diag{\operatorname{diag}}
\newcommand{\suchthat}{\;\ifnum\currentgrouptype=16 \middle\fi|\;}
\def\tr{\operatorname{tr}}
\title{Eigenvalues of locally positive semidefinite matrices: Non-convexity and Geometry}
\author{Jose Acevedo, Grigoriy Blekherman, Sebastian Debus, Seokbin Lee, Cordian Riener}
\date{Source: arXiv:2608.15444v2 (CC BY 4.0)}
\begin{document}
\maketitle

\noindent\emph{Source and licence.} Eigenvalues of locally positive semidefinite matrices: Non-convexity and Geometry. Version arXiv:2608.15444v2 (CC BY 4.0), distributed by arXiv under the Creative Commons Attribution 4.0 International licence, http://creativecommons.org/licenses/by/4.0/, as recorded in the arXiv licence metadata for this exact version and shown on the abstract page https://arxiv.org/abs/2608.15444v2. Full licence notice: \texttt{licenses/arxiv-2608.15444-CC-BY-4.0.txt}.\par\smallskip

\noindent\textit{Editorial scope.} Counted unit: the article's proposition equivalences (source0.tex lines 484--494). The article's complete original proof of it is delivered with it (source0.tex lines 898--952, one \texttt{proof} environment of 55 lines, which the article places away from the statement). No mathematics is written, repaired or re-derived: the statement and the proof are the article's own lines, in the article's own notation, and no retained line is substituted or re-flowed.  Retained uncounted as the source-local context the proof needs: lines 461--463.  Nothing was omitted from any retained span, so the package carries no omission record.  No retained line contains a citation, so the wrapper carries no bibliography and no bibliography entry.  No retained line contains a figure environment, an image-inclusion command or drawing code, so no image is delivered and the package declares no asset dependency.  Editorial formatting only, in addition: an AMS \texttt{amsart} preamble that restates the article's own declarations and macros used by the retained lines, namely the environments \texttt{lemma}, \texttt{proposition} on separate counters, together with the standard packages those lines need.  The retained environments print in this document as Lemma 1, Proposition 1 in order of appearance, whereas the article prints the same items with the numbering of its own full document.  The complete native source of the article as distributed in the arXiv e-print for this exact version is bundled unchanged as \texttt{sources/207748/source0.tex}.
\begin{lemma}\label{lem: s,t <= 0}
Minimizing $s+t$ over $(s,t) \in \mathcal{L}_{n,d}$ with $s,t\leq 0$ yields the same value as minimizing $s+t$ over the whole $\mathcal{L}_{n,d}$. 
\end{lemma}

\begin{proposition}\label{prop: equivalences}
For $2 \leq d \leq n$, consider the following two optimization problems over $n$ complex numbers $z_1,\dots,z_n$:
    \begin{align*}
    \mu_{n,d}^{\max}  & :=   \min\left\{\max_{S\in\binom{[n]}d}\sum_{s\in S}|z_s|+\left|\sum_{s\in S} z_s\right|\,\suchthat\,z_1,\dots,z_n\in\CC,\, \sum_{k=1}^n|z_k|=1\right\}, \\
    \pi_{n,d}^{\max} & :=  
    \min\left\{\max_{S\in\binom{[n]}d}\sum_{s\in S}|z_s|+\left|\sum_{s\in S} z_s\right|\,\suchthat\,z_1,\dots,z_n\in\CC,\, \sum_{k=1}^n|z_k|=1,\, \sum_{k=1}^n z_k=0\right\}.    
    \end{align*} 
    The following relations hold between $\mu_{n,d}^{\tr}$ and $ \mu_{n,d}^{\max}$, as well as between $\pi_{n,d}^{\tr}$ and $ \pi_{n,d}^{\max}$: 
    \[ \mu_{n,d}^{\tr} = 2 \left( 1- \frac{1}{\mu_{n,d}^{\max}}\right) \text{ and } \pi_{n,d}^{\tr} = 1- \frac{1}{\pi_{n,d}^{\max}}.\]
    Moreover, for an optimal solution $(z_1,\dots,z_n)$ attaining $\mu_{n,d}^{\max}$ the corresponding pair $(s,t)$ of the optimization problem for $\mu_{n,d}^{\operatorname{tr}}$ satisfies $s=1-\frac{1}{\mu_{n,d}^{\max}}(1-|\sum_{j=1}^n z_j|)$ and $t=1-\frac{1}{\mu_{n,d}^{\max}}(1+|\sum_{j=1}^n z_j|)$, up to permutation of $s$ and $t$.
\end{proposition}

\begin{proof}[Proof of Proposition \ref{prop: equivalences}]
By Lemma \ref{lem: s,t <= 0} we may assume that $\max \{s,t\} \leq 0$.
Observe that both problems optimize a continuous function over a compact set. Thus, the optimal values are attained.

Let $M \in \Sym_n$ with $M \sim \diag (1,\ldots,1,s,t)$ and $\max\{s,t\} \leq 1$. This implies that 
\[
I-M \sim \diag(0,\ldots,0,1-s,1-t) \succeq 0,
\]
so there exists a matrix $B \in \RR^{2 \times n}$ with $I-M=B^\top B$. Then $M \in \S^{n,d}$ if and only if $I_d-B_S^\top B_S \succeq 0$ for all $S \in \binom{[n]}{d}$. Note that this is equivalent to $\max_{S \in \binom{[n]}{d}} \rho(B_S^\top B_S) \leq 1$.
Let $A \in \RR^{2 \times n}$ be proportional to $B$ and express its columns in polar coordinates: 
\begin{align}\label{eq: matrix A}
A=\begin{bmatrix}
a_1\cos\b_1 & a_2\cos\b_2 & \dots & a_n\cos\b_n\\
a_1\sin\b_1 & a_2\sin\b_2 & \dots & a_n\sin\b_n
\end{bmatrix}.
\end{align}
For every index set $S \in \binom{[n]}{d}$, the two largest eigenvalues of $A_S^\top A_S$ coincide with the two eigenvalues of $A_SA_S^\top$. 
To compute the eigenvalues of $A_SA_S^\top$, we calculate its characteristic polynomial $t^2-\tr(A_SA_S^\top)t+\det(A_SA_S^\top)$, where $\tr(A_SA_S^\top)=\sum_{j \in S}a_j^2$. The determinant can be calculated by the Cauchy-Binet formula:
\begin{align*}
&\det(A_SA_S^\top) = \sum_{1 \leq k < l \leq d} \left(\det\begin{bmatrix}
    a_{j_k} \cos\beta_{j_k} & a_{j_l} \cos\beta_{j_l} \\ a_{j_k} \sin\beta_{j_k} & a_{j_l} \sin\beta_{j_l}
\end{bmatrix} \right)^2 \\ 
&= \sum_{1 \leq k < l \leq d}\left( a_{j_k}a_{j_l}\cos\beta_{j_k}\sin\beta_{j_l}-a_{j_k}a_{j_l}\cos\beta_{j_l}\sin\beta_{j_k}\right)^2 =\sum_{1 \leq k < l \leq d} a_{j_k}^2a_{j_l}^2\sin^2(\beta_{j_l}-\beta_{j_k})~.    
\end{align*} 
Setting $r_k:=a_k^2, \theta_k:=2\b_k$, we calculate with the double-angle identities the largest eigenvalue $\lambda$:
\begin{align*}
    2\lambda&=\sum_{j \in S} a_j^2+\sqrt{\left(\sum_{j \in S} a_j^2\right)^2-4\sum_{k<l \in S}a_k^2a_l^2\sin^2(\b_k-\b_l)}\\
    &=\sum_{j \in S} a_j^2+\sqrt{\left(\sum_{j \in S} a_j^2\right)^2-2\sum_{k<l \in S}a_k^2a_l^2(1-\cos(\theta_k-\theta_l))}\\
    &=\sum_{j \in S} a_j^2+\sqrt{\sum_{j \in S} a_j^4+2\sum_{k<l \in S}a_k^2a_l^2\cos(\theta_k-\theta_l)}\\
    &=\sum_{j \in S} r_j+\sqrt{\sum_{j \in S} r_j^2+2\sum_{k<l \in S}r_kr_l\cos(\theta_k-\theta_l)}\\
    &=\sum_{j\in S}|z_j|+\left|\sum_{j\in S} z_j\right|,
\end{align*}
where $z_k:=r_k(\cos\theta_k+i\sin \theta_k)$. Thus, the maximum eigenvalue $\alpha := \max_{S \in \binom{[n]}{d}}\rho(A_S^\top A_S)$ is exactly $\max_{S \in \binom{[n]}{d}} \frac{1}{2}\sum_{j\in S}|z_j|+\left|\sum_{j\in S} z_j\right|$.
We now proceed with proving each relation separately.\\
(1) We want to verify the relation between $\mu_{n,d}^{\tr}$ and $\mu_{n,d}^{\max}$. Since we optimize a linear function over $\mathcal{L}_{n,d}$, we may assume that the optimal value $(s,t)$ is attained at the boundary of $\mathcal{L}_{n,d}$. In particular, the corresponding matrix $M \in \S^{n,d}$ must be on the boundary of $\S^{n,d}$. Thus, $\max_{S \in \binom{[n]}{d}} \rho(B_S^\top B_S)=1$ holds. We rescale $B$ by its Frobenius norm and consider $A := \frac{1}{\|B\|_F} B$. Then, we have $\|A\|_F^2 =1$ and $ 1 = \max_{S \in \binom{[n]}{d}} \rho(B_S^\top B_S) = \|B\|_F^2 \alpha.$ Moreover,
    \begin{align}\label{eq: 2-s-t = 1/alpha}
    2-(s + t) = \tr (I-M) = \tr (B^\top B)  = \|B\|_{F}^2 = \frac{1}{\alpha},
    \end{align}
    so minimizing $s+t$ is equivalent to minimizing $\alpha$. Since 
    \[
    \|A\|_F^2 = \tr(AA^\top) = \sum_{k=1}^n a_k^2 = \sum_{k=1}^n |z_k|=1,
    \]
    we observe that our variables exactly parametrize the feasible region of the optimization problem $\mu_{n,d}^{\max}$, which proves $\alpha = \frac{1}{2} \mu_{n,d}^{\max}$ and together with \eqref{eq: 2-s-t = 1/alpha} we obtain \[\mu_{n,d}^{\tr}=2\left(1-\frac{1}{\mu_{n,d}^{\max}}\right).\]
    Let $z_1,\dots,z_n$ be as above and let $(s,t) \in \mathcal{L}_d$ be a corresponding optimal solution attaining $\mu_{n,d}^{\tr}$. Then $\diag(0,\dots,0,1-s,1-t)=\frac{2}{\mu_{n,d}^{\max}}A^\top A$. The two largest eigenvalues of $A^\top A$ are the same as the eigenvalues of $AA^\top$ which are $\frac{1}{2}(1\pm | \sum_{j=1}^nz_j |),$
since $\sum_{j=1}^n |z_j| =1$.
Therefore, $1-s=\frac{1}{\mu_{n,d}^{\max}}(1-|\sum_{j=1}^n z_j|)$ and $1-t=\frac{1}{\mu_{n,d}^{\max}}(1+|\sum_{j=1}^n z_j|)$, up to permutation, which implies the claim.


(2) Next, we assume that $M \sim \diag(1,\ldots,1,t,t)$. The spectral decomposition of of $M$ is $M=I+(t-1)P$, where $P$ is the orthogonal projection matrix onto the $t$-eigenspace. Since $P$ has rank 2, we can factor $P=A^\top A$ for some matrix $A \in \RR^{2 \times n}$ with orthonormal rows. Then $M \in \S^{n,d}$ if and only if  $1+(t-1)\alpha \geq 0$, where $\alpha = \max_{S \in \binom{[n]}{d}}\rho(A_SA_S^\top)$. Observe that $\alpha > 0$ since at least one diagonal entry of $P$ is positive. Thus, $M \in \S^{n,d}$ if and only if $t \geq 1- \frac{1}{\alpha}$. So, minimizing $t$ is equivalent to minimizing $\alpha$. 
    As in \eqref{eq: matrix A}, we write the columns of $A$ in polar coordinates. Because $A$ has orthonormal rows, we get $AA^\top = I_2$ which is equivalent to the three conditions: $\sum_{k=1}^n a_k^2 \cos^2 \beta_k = 1, \sum_{k=1}^n a_k^2 \sin^2 \beta_k =1$, and $\sum_{k=1}^n a_k^2 \cos\beta_k \sin \beta_k = 0$. We also have $z_k = a_k^2(\cos (2\beta_k) + \mathrm{i} \sin (2\beta_k))$, so 
    the three conditions are equivalent to $\sum_{k=1}^n |z_k| =2$ and $\sum_{k=1}^n z_k = 0$. This yields
    \[ \alpha = \frac{1}{2}\max_{S\in\binom{[n]}d} \sum_{s\in S}|z_s|+\left|\sum_{s\in S} z_s\right|. \] To align this with the feasible region of the optimization problem of $\pi_{n,d}^{\max}$ we rescale to complex numbers $\sum_{k=1}^n |z_k|=1$ and this proves the relation
    \[ \pi_{n,d}^{\tr} = 1- \frac{1}{\pi_{n,d}^{\max}}.\]

\end{proof}

\end{document}
